\section{El peso del texto frente a la imagen en los generadores}
\label{sec:disc_texto_imagen}
% Extensión estimada: 2 páginas, cuatro párrafos (ficha 6.2).
% Fuentes: sec:res_vlm_controles, sec:res_gen_sin_anclaje, sec:res_gen_fuentes, sec:res_gen_ablaciones, sec:res_gen_rag,
% sec:res_gen_comparacion, sec:res_gen_pipeline, sec:res_gen_cualitativos, sec:res_gen_estado_arte, sec:met_generacion
% (instrucción de ajuste sin diagnóstico; LoRA de Mini-VLM solo en el codificador visual); rohrbach2018object,
% bai2024hallucination. Reemplaza la antigua 6.3 y la parte de PI2 de la antigua 6.1.

Los Capítulos~\ref{cp:res_clasif} y~\ref{cp:res_gen} evaluaron a los generativos en tareas distintas, pero en las dos
aparece un mismo patrón: en las condiciones evaluadas, cuando el texto de la entrada y la imagen apuntan a respuestas distintas, el generador tiende a seguir al texto. En la clasificación, los controles son compatibles con que el \emph{prompt} fije una respuesta por defecto que la
imagen desplaza muy poco
(Sección~\ref{sec:res_vlm_controles}).

En la generación de reportes, el patrón aparece en tres situaciones. Sin un diagnóstico en la instrucción, SkinVL repite en siete de cada diez reportes un mismo texto ajeno a la imagen (Sección~\ref{sec:res_gen_sin_anclaje}). En la cascada, el diagnóstico afirmado no cambia al quitar la imagen y, con un
diagnóstico equivocado, el reporte lo afirma en lugar de corregirlo (Sección~\ref{sec:res_gen_ablaciones}). Por
último, cuando recibe a la vez el diagnóstico del clasificador y el de los vecinos, el generador sigue al clasificador
aunque los vecinos acierten: en las 72 imágenes en que el clasificador se equivoca y los vecinos aciertan, recupera la clase verdadera en una de cada seis
(Sección~\ref{sec:res_gen_rag}).

En conjunto, estos resultados muestran que, bajo las condiciones evaluadas, los generadores tienden a conservar el diagnóstico proporcionado en la instrucción, incluso cuando este no coincide con la categoría verdadera. Esta tendencia sugiere una fuerte influencia de la información textual sobre la impresión diagnóstica generada, aunque no permite determinar de manera concluyente el peso relativo de cada modalidad. Además, el diagnóstico de la instrucción pesa más que el contexto recuperado, lo que sugiere una jerarquía entre las fuentes del reporte.

El origen de esa influencia no parece estar en el ajuste de este trabajo. La instrucción utilizada durante el ajuste, «Describe this skin lesion.», no proporciona explícitamente un diagnóstico que el modelo deba reproducir (Sección~\ref{sec:met_generacion}). Además, en Mini-VLM el ajuste no llega al modelo de lenguaje, y SkinVL, que no se ajustó en este trabajo,
transcribe el diagnóstico del mismo modo. Una explicación posible es que la tendencia proviene del modelo de lenguaje
de base, que en su ajuste original aprendió a seguir lo que la instrucción afirma, y que la imagen no alcanza a
contradecir. Es una forma del peso de los patrones del lenguaje que la literatura señala como causa de las
alucinaciones~\cite{rohrbach2018object,bai2024hallucination}, aunque aquí no produce un hallazgo inventado, sino
la adopción de un diagnóstico ajeno. La explicación alternativa es que ceder sea lo esperable en un generador que no
distingue las clases por sí mismo (Sección~\ref{sec:disc_representacion}): si el modelo no logra obtener de la imagen información suficiente para distinguir las categorías, puede tender a utilizar el diagnóstico proporcionado en la instrucción como principal referencia para redactar la conclusión. Las dos explicaciones predicen lo mismo en estos experimentos,
y separarlas requeriría un generador que clasifique bien sin anclaje, que ninguno de los evaluados es.

Cualquiera sea su origen, esa influencia condiciona cómo debe evaluarse un sistema que integra un clasificador y un generador. Como el generador transcribe el
diagnóstico, la fidelidad diagnóstica del reporte mejora con el clasificador más que con el generador: cambiar la
fuente de DermLIP sin ajuste a Swin-Tiny agrega 0.132 de consistencia con el mismo generador
(Sección~\ref{sec:res_gen_fuentes}), mientras que ninguno de los tres generadores eleva la consistencia más de unas milésimas por encima de la exactitud de la fuente. Evaluado como un todo, el sistema puede atribuirle al generador una fidelidad que proviene del clasificador, y por eso debe medirse por componente. En los trabajos incluidos en la revisión no se identificó una evaluación que distinguiera, mediante comparaciones equivalentes, la contribución del clasificador y del generador a la consistencia diagnóstica del reporte final (Sección~\ref{sec:res_gen_estado_arte}). La Tabla~\ref{tab:res_gen_aportes} resume esa separación para este trabajo.

En consecuencia, la selección del generador prioriza su capacidad para conservar el diagnóstico recibido y evitar errores adicionales, dado que su contribución principal al sistema es la producción del contenido textual. LLaVA-Med omite el diagnóstico recibido en uno de cada siete
reportes aun con el \emph{oracle}, y en el nivel diferencial deja los candidatos dentro de un razonamiento sin cerrar
(Secciones~\ref{sec:res_gen_fuentes} y~\ref{sec:res_gen_pipeline}). Por eso el \emph{pipeline} usa Mini-VLM, que conserva el diagnóstico como SkinVL, describe tanto
como LLaVA-Med y entra en una GPU de consumo.

En estos experimentos, esa misma influencia impide que el generador funcione como una segunda opinión, y además alcanza a la descripción. Sin la
imagen, cada generador escribe un único texto por clase, una descripción prototípica del diagnóstico que
recibe; con ella, la descripción varía entre lesiones (Sección~\ref{sec:res_gen_ablaciones}). Una lectura posible es
que el reporte parte de ese prototipo y la imagen lo modifica, pero sin contradecirlo, como en el melanoma afirmado
como nevo y descrito con una pigmentación homogénea «with no visible melanocytic infiltration or atypical features» (Sección~\ref{sec:res_gen_pipeline}).
Qué parte de la descripción proviene del prototipo y qué parte de la imagen no se midió. Lo que sí se sigue de los
resultados es que la verosimilitud de un reporte no indica si su diagnóstico es correcto, porque un error del
clasificador llega al reporte acompañado de rasgos que lo justifican. En este diseño, la señal de cuándo desconfiar debe venir, por lo tanto, de fuera del generador, y esa es la función del ruteo, tema de la sección siguiente.
